\chapter{Internal Peripheral Drivers}
\label{app:internal}

This appendix is a per-driver reference for the SoC's \emph{on-chip} peripherals
--- the blocks built into the ESP32-S3 itself (GPIO, the serial buses, the
timers and PWM units, DMA, the analog and audio paths, the low-power domain, and
the crypto accelerators). The companion appendix~\ref{app:external} covers the
drivers for \emph{external} parts wired to these peripherals. The narrative HAL
chapters (\S\ref{ch:hal} onward) introduce the conventions and the design
rationale; this appendix collects, for each driver, a short description, a digest
of its public API, and a worked example drawn from the \texttt{examples/} tree.

A few patterns recur and are worth stating once. The bus drivers (UART, SPI,
I2C, I2S, LCD, TWAI) seal their unsynchronised register access in a private
\texttt{.Engine} child and hand the application an acquired, auto-releasing
\texttt{Session}; the resource drivers (LEDC, MCPWM, Timer, PCNT, SDM, RMT, GDMA,
ADC) hand out a \emph{limited, controlled} handle that claims hardware on
\texttt{Claim}/\texttt{Acquire} and frees it on scope exit, even during exception
unwinding. Both idioms depend on finalization, so those drivers require the
\texttt{embedded} or \texttt{full} profile; the register-poke drivers (GPIO,
RNG, RTC, RTC-IO, Touch, SHA, AES, SDMMC) work under every profile, and each
section notes which case applies. See the driver-design chapter
(\S\ref{ch:driver-design}) for why the ownership gateway is built this way.

\section{GPIO --- General-Purpose I/O}\index{GPIO driver}
\label{app:int:gpio}

The ESP32-S3 routes every digital pad through a \emph{GPIO matrix}: a pad can be
a software-driven GPIO or be cross-connected to any peripheral's input/output
signal. The Ada driver \texttt{ESP32S3.GPIO} sits on top of the generated
register layer (\texttt{Interfaces.ESP32S3.GPIO} plus \texttt{IO\_MUX}) and
exposes the plain-GPIO case: configure a pad's direction, pull and drive
strength, then \texttt{Set} / \texttt{Clear} / \texttt{Toggle} / \texttt{Write}
/ \texttt{Read} it. \texttt{Configure} always selects the matrix path
(\texttt{MCU\_SEL = 1}, output index 256); routing a pad to a \emph{peripheral}
signal instead is the job of that peripheral's own \texttt{Configure\_Pins}.

The pad space is modelled by \texttt{Pad\_Number} (\texttt{-1\,..\,48}), with
\texttt{-1} the \texttt{No\_Pin} sentinel for an optional unrouted line. The
subtype \texttt{Pin\_Id} carries a \emph{static} predicate that excludes the
pads the silicon does not bring out (22\,..\,25) and those bonded to the
in-package flash and octal PSRAM (26\,..\,37): naming one of those as a static
value is a compile-time error, and a dynamic value is caught at run time under
\texttt{-gnata}. \texttt{Set} / \texttt{Clear} / \texttt{Write} use the
hardware-atomic W1TS/W1TC banks and \texttt{Read} is a pure load, so those are
concurrency-safe as written; the read-modify-write paths (\texttt{Configure},
\texttt{Toggle}) are serialised through a protected object, so the package needs
a tasking runtime. A companion child, \texttt{ESP32S3.GPIO.Interrupts}, owns the
single GPIO interrupt source, demultiplexes it by latched status, and dispatches
a per-pin \texttt{Callback} that runs in interrupt context under the level-3
ceiling.

\begin{lstlisting}
package ESP32S3.GPIO is
   type Pad_Number is range -1 .. 48;
   No_Pin : constant Pad_Number := -1;

   subtype Pin_Id is Pad_Number range 0 .. 48 with
     Static_Predicate => Pin_Id in 0 .. 21 | 38 .. 48;
   subtype Optional_Pin is Pad_Number with
     Static_Predicate => Optional_Pin in -1 | 0 .. 21 | 38 .. 48;

   type Pin_Mode       is (Input, Output);
   type Pull_Mode      is (Floating, Pull_Up, Pull_Down);
   type Drive_Strength is (Drive_Weak, Drive_Medium, Drive_Strong, Drive_Strongest);

   procedure Configure (Pin : Pin_Id; Mode : Pin_Mode;
                        Pull  : Pull_Mode      := Floating;
                        Drive : Drive_Strength := Drive_Medium);
   procedure Set    (Pin : Pin_Id);   --  atomic W1TS
   procedure Clear  (Pin : Pin_Id);   --  atomic W1TC
   procedure Toggle (Pin : Pin_Id);
   procedure Write  (Pin : Pin_Id; On : Boolean);
   function  Read   (Pin : Pin_Id) return Boolean;
end ESP32S3.GPIO;

--  Pin interrupts (ESP32S3.GPIO.Interrupts):
type Trigger  is (Rising_Edge, Falling_Edge, Any_Edge, Low_Level, High_Level);
type Callback is access procedure;
procedure Enable  (Pin : Pin_Id; On : Trigger; Action : Callback);
procedure Disable (Pin : Pin_Id);
\end{lstlisting}

The \texttt{gpio0\_blink} example drives GPIO0 from a library-level task ---
configure once as a push-pull output, then toggle on a fixed period:

\begin{lstlisting}
with ESP32S3.GPIO;
with Ada.Real_Time; use Ada.Real_Time;

Pin : constant ESP32S3.GPIO.Pin_Id := 0;
task Blinker;
task body Blinker is
   Period : constant Time_Span := Milliseconds (250);   --  ~2 Hz
   Next   : Time;
   High   : Boolean := False;
begin
   ESP32S3.GPIO.Configure
     (Pin, ESP32S3.GPIO.Output, Drive => ESP32S3.GPIO.Drive_Strong);
   Next := Clock + Period;
   loop
      delay until Next;
      High := not High;
      ESP32S3.GPIO.Write (Pin, High);
      Next := Next + Period;
   end loop;
end Blinker;
\end{lstlisting}

\cnote{ESP-IDF wraps the same matrix in \texttt{gpio\_config()} +
\texttt{gpio\_set\_level()}, taking a raw \texttt{int} pin number checked only at
run time. The Ada \texttt{Pin\_Id} predicate moves that check to compile time for
the common static-pin case --- a flash/PSRAM pad never even links.}

\noindent\textbf{Example:} \texttt{examples/esp32s3\_gpio0\_blink}.

\section{UART --- Asynchronous Serial}\index{UART driver}
\label{app:int:uart}

The chip has three general-purpose UART controllers; the bare runtime uses
USB-Serial-JTAG for its console, so even UART0's pads are free to repurpose.
\texttt{ESP32S3.UART} is the only interface the application sees: the raw
register driver lives in a private child, \texttt{ESP32S3.UART.Engine}, that
cannot be \texttt{with}-ed from outside the subtree, so the unsynchronised
primitives can never be reached by accident.

Configuration is split in two. \texttt{Setup} runs once per port at startup
(single-threaded) and fixes baud, frame format and pin routing --- all four
lines (\texttt{Tx}, \texttt{Rx}, \texttt{Rts}, \texttt{Cts}) are
\texttt{Optional\_Pin}, so a link routes only what it uses. Everything after
that goes through an acquired \texttt{Session}: \texttt{Acquire} hands out a
limited, non-copyable token that owns the port exclusively (other tasks suspend
until it is released), and every transfer \emph{and} every later
reconfiguration (\texttt{Set\_Baud}, \texttt{Configure\_Pins},
\texttt{Set\_Inversion}, \ldots) is a method on the held \texttt{Session} --- so
a reconfigure can never race another task's transfer. The \texttt{Session} is a
controlled type that auto-releases on scope exit (including during exception
unwinding), so this driver, like the other task-safe HAL drivers, targets the
\texttt{embedded}/\texttt{full} profile (full exceptions and finalization).
\texttt{Enable\_Loopback} ties TX back to RX inside the controller for a
no-wiring self-test.

\begin{lstlisting}
package ESP32S3.UART is
   type UART_Port  is (UART0, UART1, UART2);
   subtype Baud_Rate is Positive range 300 .. 5_000_000;
   type Data_Bits   is range 5 .. 8;
   type Parity_Mode is (None, Even, Odd);
   type Stop_Bits   is (One, Two);
   type Byte        is new Interfaces.Unsigned_8;
   type Byte_Array  is array (Natural range <>) of Byte;
   type Session     is limited private;            --  auto-releasing

   procedure Setup (Port : UART_Port;
                    Baud : Baud_Rate   := 115_200;
                    Bits : Data_Bits   := 8;
                    Parity : Parity_Mode := None;
                    Stop   : Stop_Bits   := One;
                    Tx, Rx, Rts, Cts : ESP32S3.GPIO.Optional_Pin
                                       := ESP32S3.GPIO.No_Pin;
                    Rx_Flow_Threshold : Natural := 100);

   procedure Acquire (S : in out Session; Port : UART_Port);
   procedure Write   (S : Session; Data : Byte_Array);
   procedure Read    (S : Session; Data : out Byte_Array; Count : out Natural);
   function  Available (S : Session) return Natural;
   procedure Enable_Loopback (S : Session; On : Boolean := True);
   procedure Release (S : in out Session);
   Not_Initialized, Not_Owned : exception;
end ESP32S3.UART;
\end{lstlisting}

The \texttt{uart\_loopback} self-test brings up UART1, ties TX to RX internally,
and round-trips a known buffer --- no wiring needed:

\begin{lstlisting}
with ESP32S3.UART; use ESP32S3.UART;

Port : constant UART_Port := UART1;
Tx   : constant Byte_Array := (16#55#, 16#AA#, 16#00#, 16#FF#, 16#12#, 16#34#);
S    : Session;
Rx   : Byte_Array (Tx'Range);
Got  : Natural;
begin
   Setup (Port, Baud => 115_200);     --  8-N-1 defaults
   Acquire (S, Port);
   Enable_Loopback (S);               --  internal TX->RX
   Write (S, Tx);
   Read  (S, Rx, Got);
   Release (S);
   --  Got = Tx'Length and Rx = Tx  on a healthy data path
\end{lstlisting}

\noindent\textbf{Example:} \texttt{examples/esp32s3\_uart\_loopback}.

\section{SPI --- Serial Peripheral Interface}\index{SPI driver}
\label{app:int:spi}

\texttt{ESP32S3.SPI} is a task-safe, full-duplex \emph{master} over the two
general-purpose hosts, \texttt{SPI2} and \texttt{SPI3} (\texttt{SPI0}/\texttt{SPI1}
drive the in-package flash and PSRAM and are deliberately not offered). As with
UART, the unsynchronised register driver is sealed in a private child,
\texttt{ESP32S3.SPI.Engine}, and reached only through an acquired
\texttt{Session}.

\texttt{Setup} brings a host up at one of the four clock modes (CPOL/CPHA 0\,..\,3)
and a bit clock from roughly 80\,kHz to 80\,MHz, and claims a GDMA channel;
\texttt{Set\_Clock} re-tunes the rate without re-claiming DMA, which lets a
device driver init slowly then run fast (an SD card negotiates at
$\le$400\,kHz, then switches to several MHz). \texttt{Configure\_Pins} routes
\texttt{Sclk}/\texttt{Mosi}/\texttt{Miso}/\texttt{Cs} to pads (each an
\texttt{Optional\_Pin}), and \texttt{Enable\_Loopback} ties MOSI to MISO through
one pad for a self-test. The single transfer primitive is a blocking,
full-duplex DMA \texttt{Transfer} of 1\,..\,4095 bytes between two SRAM buffers,
addressed by \texttt{System.Address}. A device that brings its own chip-select ---
to share the bus, or behind a 3:8 decoder or an I/O-expander --- registers a
library-level \texttt{CS\_Select} callback at \texttt{Acquire} and brackets its
transaction with \texttt{Select\_Device}; the driver then suppresses the hardware
\texttt{CS0} for that session. The controlled \texttt{Session} auto-releases on
scope exit, so this driver targets \texttt{embedded}/\texttt{full}.

\begin{lstlisting}
package ESP32S3.SPI is
   type SPI_Host is (SPI2, SPI3);
   subtype SPI_Mode is Natural range 0 .. 3;
   type Session is limited private;            --  auto-releasing

   procedure Setup       (Host : SPI_Host; Mode : SPI_Mode := 0;
                          Clock_Hz : Positive := 1_000_000);
   procedure Set_Clock   (Host : SPI_Host; Hz : Positive);
   procedure Enable_Loopback (Host : SPI_Host; Pad : ESP32S3.GPIO.Pin_Id);
   type CS_Select is access procedure            --  app chip select: library-
     (Ctx : System.Address; Active : Boolean);   --  level, closure-free

   procedure Configure_Pins (Host : SPI_Host;
                             Sclk, Mosi, Miso : ESP32S3.GPIO.Optional_Pin;
                             Cs   : ESP32S3.GPIO.Optional_Pin := No_Pin);

   procedure Acquire  (S : in out Session; Host : SPI_Host;
                       Select_CB : CS_Select := null;
                       Ctx : System.Address := System.Null_Address);
   procedure Select_Device (S : in out Session; On : Boolean);
   procedure Transfer (S : Session; Tx, Rx : System.Address; Length : Natural);
   procedure Release  (S : in out Session);
   Not_Initialized, Not_Owned : exception;
end ESP32S3.SPI;
\end{lstlisting}

A no-wiring round-trip: bring up SPI2, loop MOSI to MISO on a free pad, and DMA
a buffer through it --- what comes back on \texttt{Rx} equals what went out:

\begin{lstlisting}
with ESP32S3.SPI; use ESP32S3.SPI;
with Interfaces;

Tx : aliased array (0 .. 7) of Interfaces.Unsigned_8
       := (16#01#, 16#23#, 16#45#, 16#67#, 16#89#, 16#AB#, 16#CD#, 16#EF#);
Rx : aliased array (Tx'Range) of Interfaces.Unsigned_8 := (others => 0);
S  : Session;
begin
   Setup (SPI2, Mode => 0, Clock_Hz => 1_000_000);
   Enable_Loopback (SPI2, Pad => 8);          --  MOSI -> MISO on GPIO8
   Acquire (S, SPI2);
   Transfer (S, Tx'Address, Rx'Address, Tx'Length);
   Release (S);                                --  Rx now mirrors Tx
\end{lstlisting}

The same driver carries real traffic in \texttt{esp32s3\_sd\_spi} (an SD card on
\texttt{SPI2}, init at 400\,kHz then data at 8\,MHz) and behind the ST7789
display driver.

\noindent\textbf{Example:} \texttt{examples/esp32s3\_sd\_spi} (and the ST7789
display) exercise this driver against real hardware.

\section{I2C --- Two-Wire Master}\index{I2C driver}
\label{app:int:i2c}

\texttt{ESP32S3.I2C} is a polled, single-segment master over the two
controllers \texttt{I2C0} and \texttt{I2C1}. Each \texttt{Write} or
\texttt{Read} is one complete \texttt{START}\,\ldots\,\texttt{STOP} transaction
of up to \texttt{Max\_Transfer} (32) bytes --- the controller's FIFO depth.
\texttt{Setup} fixes the SCL bit clock (100\,kHz standard / 400\,kHz fast are
typical), and \texttt{Configure\_Pins} routes \texttt{Scl}/\texttt{Sda} as
open-drain lines with the internal pull-ups on, enough for bring-up (add real
resistors for a production bus). The same sealed-\texttt{Engine},
\texttt{Acquire}-a-\texttt{Session} structure as the other buses applies, and the
controlled \texttt{Session} auto-releases, so the package targets
\texttt{embedded}/\texttt{full}.

A 7-bit \texttt{Slave\_Address} is given without the R/W bit (the driver adds
it). \texttt{Write} reports \texttt{Success} = \texttt{True} only if the slave
ACKed the address and every byte; a zero-length \texttt{Write} is an
address-only probe, which makes bus scanning trivial, and \texttt{Check\_Ack =>
False} clocks the frame out regardless of ACK. Unlike UART/SPI there is no
on-chip loopback: I2C's SDA is a bidirectional wired-AND node and the matrix
gives each pad one output source, so verifying the read/ACK direction needs a
real device or a jumper.

\begin{lstlisting}
package ESP32S3.I2C is
   type I2C_Host is (I2C0, I2C1);
   subtype Slave_Address is Natural range 0 .. 16#7F#;
   type Byte       is new Interfaces.Unsigned_8;
   type Byte_Array is array (Natural range <>) of Byte;
   Max_Transfer : constant := 32;
   type Session  is limited private;            --  auto-releasing

   procedure Setup          (Host : I2C_Host; Clock_Hz : Positive := 100_000);
   procedure Configure_Pins (Host : I2C_Host;
                             Scl, Sda : ESP32S3.GPIO.Pin_Id);

   procedure Acquire (S : in out Session; Host : I2C_Host);
   procedure Write   (S : Session; Addr : Slave_Address; Data : Byte_Array;
                      Success : out Boolean; Check_Ack : Boolean := True);
   procedure Read    (S : Session; Addr : Slave_Address; Data : out Byte_Array;
                      Success : out Boolean);
   procedure Release (S : in out Session);
   Not_Initialized, Not_Owned : exception;
end ESP32S3.I2C;
\end{lstlisting}

The \texttt{i2c\_loopback} self-test brings up \texttt{I2C0} and probes an empty
bus: an ACK-checked write to an absent address must come back NACKed
(\texttt{Success = False}), proving START, addressing and ACK detection on
silicon:

\begin{lstlisting}
with ESP32S3.I2C; use ESP32S3.I2C;

Host   : constant I2C_Host      := I2C0;
Absent : constant Slave_Address := 16#55#;   --  nothing answers here
S      : Session;
Ok     : Boolean;
begin
   Setup (Host, Clock_Hz => 100_000);
   Configure_Pins (Host, Scl => 6, Sda => 4);
   Acquire (S, Host);
   Write (S, Absent, (1 => 16#00#), Ok);     --  address-only-ish probe
   Release (S);
   --  PASS when (not Ok): the master saw the NACK from the empty bus
\end{lstlisting}

\noindent\textbf{Example:} \texttt{examples/esp32s3\_i2c\_loopback}.

\section{RNG --- Hardware Random Number Generator}\index{RNG driver}
\label{app:int:rng}

\texttt{ESP32S3.RNG} reads the on-chip random number generator: each load of its
data register yields a fresh 32-bit word, drawn from analog and clock-jitter
noise. The package is deliberately the simplest in the HAL --- \texttt{Read} is a
single atomic register load and \texttt{Fill} writes only the caller's buffer,
so there is no shared mutable state and \emph{no} protected object. It is marked
\texttt{Preelaborate} and stays usable from a ZFP (no-runtime) context as well
as the tasking profiles; concurrent readers simply get independent words.

Two caveats carry over from the spec. First, this is the RNG \emph{peripheral}
register, not ESP-IDF's \texttt{WDEV\_RND\_REG} --- that one only yields entropy
with the RF clock domain up (Wi-Fi/BT), which the bare runtime never starts, so
it would read a constant here. Second, without an active hardware entropy source
(RF, or SAR-ADC bootloader entropy) the output is fine for dithering, non-secret
IDs, test data or seeding a software PRNG, but should \emph{not} be treated as a
CSPRNG. Don't poll faster than the hardware refreshes, or successive words may
correlate.

\begin{lstlisting}
package ESP32S3.RNG with Preelaborate is
   subtype Word is ESP32S3_Registers.UInt32;
   function Read return Word with Inline;       --  one fresh 32-bit value

   type Byte       is mod 2 ** 8 with Size => 8;
   type Byte_Array is array (Natural range <>) of Byte with Pack;
   procedure Fill (Buffer : out Byte_Array);    --  any length, word at a time
end ESP32S3.RNG;
\end{lstlisting}

Reading a word and filling an arbitrary-length buffer:

\begin{lstlisting}
with ESP32S3.RNG;

W   : ESP32S3.RNG.Word;
Key : ESP32S3.RNG.Byte_Array (0 .. 11);    --  12 random bytes (3 words + 0 tail)
begin
   W := ESP32S3.RNG.Read;        --  e.g. seed a software PRNG with W
   ESP32S3.RNG.Fill (Key);       --  fill the whole buffer with random bytes
\end{lstlisting}

\noindent\textbf{Example:} no dedicated example; the driver is a one-call read
(\texttt{ESP32S3.RNG.Read}) usable directly from any profile, ZFP included.

\section{Temperature --- On-Chip Die Sensor}\index{Temperature driver}
\label{app:int:temperature}

\texttt{ESP32S3.Temperature} reads the SoC's internal temperature sensor. It
reports \emph{die} temperature, not ambient air: the chip self-heats under load,
so an idle board reads a few degrees above room temperature, with roughly
$\pm$1\,\textdegree C accuracy after the factory trim and best near the middle of
the selected range. Bring-up follows ESP-IDF's \texttt{temperature\_sensor\_ll}
sequence --- gate and pulse-reset the SAR clock, open the analog REGI2C bus,
program the DAC range (a ROM call), then power the sensor up.

\texttt{Initialize} is optional: the first \texttt{Read\_*} brings the sensor up
with the default range if you skip it, and you may call \texttt{Initialize}
again to switch range. Five \texttt{Measure\_Range} values bracket different
spans; pick the one that contains your expected die temperature for best
accuracy. Readings come as signed centi-degrees
(\texttt{Read\_Centi\_Celsius}, e.g.\ 1807 = 18.07\,\textdegree C), whole degrees
(\texttt{Read\_Celsius}), or the raw 8-bit code (\texttt{Read\_Raw}) if you
prefer your own curve. The sensor is owned by a protected object so concurrent
reads are serialised (each busy-waits microseconds for a fresh conversion under
the lock), which means the package requires a tasking runtime.

\begin{lstlisting}
package ESP32S3.Temperature is
   type Measure_Range is
     (Range_Minus10_80,    --  -10 ..  80 C  (default)
      Range_20_100,        --   20 .. 100 C
      Range_50_125,        --   50 .. 125 C
      Range_Minus30_50,    --  -30 ..  50 C
      Range_Minus40_20);   --  -40 ..  20 C

   procedure Initialize (Span : Measure_Range := Range_Minus10_80);
   function  Read_Centi_Celsius return Integer;    --  1807 = 18.07 C
   function  Read_Celsius       return Integer;    --  truncated toward zero
   function  Read_Raw           return ESP32S3_Registers.Byte;
end ESP32S3.Temperature;
\end{lstlisting}

Initialise for a room-temperature board, then read the die temperature:

\begin{lstlisting}
with ESP32S3.Temperature;

Centi : Integer;
Whole : Integer;
begin
   ESP32S3.Temperature.Initialize (ESP32S3.Temperature.Range_Minus10_80);
   Centi := ESP32S3.Temperature.Read_Centi_Celsius;   --  e.g. 4231 = 42.31 C
   Whole := ESP32S3.Temperature.Read_Celsius;         --  e.g. 42
\end{lstlisting}

\noindent\textbf{Example:} no dedicated example; the driver is a direct read
(\texttt{ESP32S3.Temperature.Read\_Celsius}) after the optional
\texttt{Initialize}.

\section{LEDC --- LED Controller (PWM)}\index{LEDC driver}
\label{app:int:ledc}

The ESP32-S3's LEDC block is the simple ``dim an LED / generate a clean PWM''
peripheral: eight low-speed channels fed by four timers. A channel selects a
timer (which fixes its frequency and duty resolution) and drives a GPIO with a
duty cycle changeable at run time. Unlike MCPWM, it has no dead-time, fault or
capture logic.

The Ada driver in \texttt{ESP32S3.LEDC} hands a channel out as a \texttt{Channel}
handle that is \emph{limited} (non-copyable, so two tasks cannot drive one
channel) and \emph{controlled} (it stops the output and returns the channel to
the free pool on scope exit, including on an exception). Because you exclusively
own a claimed channel, \texttt{Set\_Duty} needs no lock. The finalization
dependency targets the embedded/full profile, not the light-tasking build.

One chip caveat surfaces in the API: a channel uses timer \texttt{(Index mod 4)},
so channels whose indices differ by 4 (e.g.\ 0 and 4) share one frequency and
resolution. Use channels 0..3 for four independent frequencies. Achievable
frequency depends on the duty bits: $f_{\max} = 80\,\mathrm{MHz} / 2^{\texttt{Bits}}$.

\begin{lstlisting}
type Channel_Index is range 0 .. 7;
subtype Resolution    is Positive range 1 .. 14;   --  duty-cycle bits
subtype Duty_Percent  is Float range 0.0 .. 100.0;

type Channel is limited private;                   --  non-copyable, auto-releasing

procedure Claim     (C : in out Channel; Index : Channel_Index);
function  Is_Valid  (C : Channel) return Boolean;
procedure Release   (C : in out Channel);
procedure Configure (C    : in out Channel;
                     Freq : Positive;
                     Pin  : ESP32S3.GPIO.Pin_Id;
                     Bits : Resolution := 10);
procedure Set_Duty  (C : Channel; Percent : Duty_Percent);  --  loads next period
procedure Stop      (C : Channel);                          --  force output low
\end{lstlisting}

\begin{lstlisting}
declare
   Ch0    : Channel;
   Duties : constant array (1 .. 2) of Duty_Percent := (25.0, 75.0);
begin
   Claim (Ch0, 0);
   Configure (Ch0, Freq => 5_000, Pin => 4, Bits => 10);   --  5 kHz, 10-bit
   for I in Duties'Range loop
      Set_Duty (Ch0, Duties (I));
      delay until Clock + Milliseconds (5);                --  let the duty latch
      --  (the self-test now GPIO-samples pin 4 to verify duty + frequency)
   end loop;
end;   --  Ch0 finalizes: output stopped, channel freed
\end{lstlisting}

\noindent\textbf{Example:} \texttt{examples/esp32s3\_ledc\_pwm}.

\section{MCPWM --- Motor-Control PWM}\index{MCPWM driver}
\label{app:int:mcpwm}

MCPWM is the heavyweight PWM block: two units (\texttt{MCPWM0}, \texttt{MCPWM1}),
each with three generator channels and three capture channels. A generator
produces edge-aligned PWM on output A (high at period start, low when the
up-counter reaches the duty comparator). It can also drive a \emph{complementary}
pair --- an inverted B output with programmable dead-time inserted on each edge so
A and B are never high together (the half-bridge / H-bridge motor case) --- plus
carrier (chopper) modulation, fault/trip-zone shutdown, and edge capture.

The Ada driver in \texttt{ESP32S3.MCPWM} models both a generator channel
(\texttt{Channel}) and a capture channel (\texttt{Capture}) as limited, controlled
handles. A unit's clock must be brought up once with \texttt{Setup} (PWM clock =
160\,MHz) before any of its channels can be claimed; claiming a channel of an
un-set-up unit raises \texttt{Not\_Initialized}. Capture timestamps are in
\texttt{Capture\_Clock\_Hz} (80\,MHz) ticks. A claimed handle's operations,
including \texttt{Set\_Duty}, need no further locking.

\cnote{ESP-IDF's \texttt{mcpwm\_*} component spreads a generator across timer,
operator, comparator and generator objects you wire together. Here one
\texttt{Configure\_Channel} call (frequency, pin, optional complement pin,
dead-time) does the equivalent, and ownership is enforced by the type system
rather than by convention.}

\begin{lstlisting}
type MCPWM_Unit    is (MCPWM0, MCPWM1);
type Channel_Index is (Ch0, Ch1, Ch2);
type Channel is limited private;     type Capture is limited private;
Not_Initialized : exception;

procedure Setup   (Unit : MCPWM_Unit);                       --  once per unit
procedure Claim   (C : in out Channel; Unit : MCPWM_Unit; Index : Channel_Index);
procedure Configure_Channel
  (C : Channel; Freq : Positive; Pin : ESP32S3.GPIO.Pin_Id;
   Complement_Pin : ESP32S3.GPIO.Optional_Pin := ESP32S3.GPIO.No_Pin;
   Dead_Time_Ns   : Natural := 0);
procedure Start (C : Channel);   procedure Stop (C : Channel);
procedure Set_Duty (C : Channel; Percent : Duty_Percent);
procedure Set_Carrier (C : Channel; Enable : Boolean := True; ...);
--  Fault / trip-zone:
procedure Configure_Fault (Unit : MCPWM_Unit; Input : Fault_Input;
                           Pin : ESP32S3.GPIO.Pin_Id; Active_High : Boolean := True);
procedure Protect_Channel (C : Channel; Input : Fault_Input;
                           Mode : Fault_Mode := One_Shot; Action : Trip_Action := Force_Low);
procedure Clear_Fault (C : Channel);   function Faulted (C : Channel) return Boolean;
--  Capture:
Capture_Clock_Hz : constant := 80_000_000;
procedure Configure_Capture (Cap : Capture; Pin : ESP32S3.GPIO.Pin_Id;
                             Edge : Cap_Edge := Both_Edges);
function  Capture_Pending (Cap : Capture) return Boolean;
procedure Read_Capture (Cap : Capture; Value : out Interfaces.Unsigned_32;
                        Falling : out Boolean);
\end{lstlisting}

\begin{lstlisting}
Setup (MCPWM0);
Claim (Gen0, MCPWM0, Ch0);
Configure_Channel (Gen0, Freq => 20_000, Pin => 4);      --  20 kHz, edge-aligned
Start (Gen0);
Set_Duty (Gen0, 75.0);

--  Complementary half-bridge pair with 1 us dead-time on channel 1:
Claim (Gen1, MCPWM0, Ch1);
Configure_Channel (Gen1, Freq => 20_000, Pin => 6,
                   Complement_Pin => 7, Dead_Time_Ns => 1_000);
Start (Gen1);  Set_Duty (Gen1, 50.0);   --  A and B each ~50 %, never overlapping

--  Over-current trip: force channel 0 low when Fault0 asserts, latched:
Configure_Fault (MCPWM0, Fault0, Pin => 8, Active_High => True);
Protect_Channel (Gen0, Fault0, One_Shot, Force_Low);
--  ... later, once the fault clears:  Clear_Fault (Gen0);
\end{lstlisting}

\noindent\textbf{Example:} \texttt{examples/esp32s3\_mcpwm\_pwm}.

\section{Timer --- General-Purpose Timers}\index{Timer driver}
\label{app:int:timer}

The S3's two timer groups, \texttt{TIMG0} and \texttt{TIMG1}, each contain one
54-bit up/down counter clocked from the APB clock through a 16-bit prescaler,
with a programmable alarm. (The per-group watchdogs are a separate block and are
left untouched.) The Ada driver in \texttt{ESP32S3.Timer} exposes one
general-purpose timer per group as a limited, controlled \texttt{Timer} handle,
released --- and the counter stopped --- automatically on scope exit.

\texttt{Configure} sets the tick rate in ticks per second (default 1\,MHz =
1\,\textmu{}s per tick; the range is roughly \texttt{80\_000\_000} down to
\texttt{80\_000\_000/65536}) and leaves the counter stopped at 0.
\texttt{Set\_Alarm} arms a one-shot: \texttt{Alarm\_Fired} reports the latched
interrupt-status flag, which stays set until \texttt{Clear\_Alarm}. Counts are
read back as the 64-bit \texttt{Ticks} type. As a finalization-using driver it
targets the embedded/full profile.

\begin{lstlisting}
type Timer_Index is range 0 .. 1;       --  0 = TIMG0, 1 = TIMG1
type Ticks       is new Interfaces.Unsigned_64;
type Timer is limited private;

procedure Claim     (T : in out Timer; Index : Timer_Index);
function  Is_Valid  (T : Timer) return Boolean;
procedure Release   (T : in out Timer);
procedure Configure (T : in out Timer; Tick_Hz : Positive := 1_000_000);
procedure Start (T : Timer);   procedure Stop (T : Timer);   procedure Reset (T : Timer);
function  Value (T : Timer) return Ticks;
procedure Set_Alarm  (T : Timer; At_Ticks : Ticks);
function  Alarm_Fired (T : Timer) return Boolean;
procedure Clear_Alarm (T : Timer);
\end{lstlisting}

\begin{lstlisting}
declare
   T : Timer;
begin
   Claim (T, 0);
   Configure (T, Tick_Hz => 1_000_000);          --  1 tick = 1 us
   Reset (T);  Start (T);
   delay until Clock + Milliseconds (50);
   pragma Assert (Value (T) in 49_000 .. 51_000); --  ~50_000 ticks, cross-checked

   Reset (T);
   Set_Alarm (T, 30_000);                         --  fire at 30 ms
   Start (T);
   while not Alarm_Fired (T) loop null; end loop;
   Clear_Alarm (T);  Stop (T);
end;   --  T finalizes: stopped and released
\end{lstlisting}

\noindent\textbf{Example:} \texttt{examples/esp32s3\_timer\_count}.

\section{PCNT --- Pulse Counter}\index{PCNT driver}
\label{app:int:pcnt}

PCNT counts edges on an input signal. The S3 has four counter units, each
accumulating into a \emph{signed 16-bit} counter that wraps at $\pm 32768$; the
classic uses are quadrature decoding, tachometers and flow-meters. The Ada driver
in \texttt{ESP32S3.PCNT} exposes the common ``count edges on a pin'' case --- the
per-unit direction-control input and the threshold-event comparators are left at
their pass-through defaults.

A unit is handed out as a limited, controlled \texttt{Unit} handle (auto-released
on scope exit). \texttt{Configure} routes a pin into the counter and starts it; by
default each rising edge increments, or set \texttt{Both\_Edges} to count both.
\texttt{Count} returns the signed value as an \texttt{Integer}, and
\texttt{Pause}/\texttt{Resume} hold and continue counting while retaining the
count.

\begin{lstlisting}
type Unit_Index is range 0 .. 3;        --  the four counter units
type Unit is limited private;

procedure Claim     (U : in out Unit; Index : Unit_Index);
function  Is_Valid  (U : Unit) return Boolean;
procedure Release   (U : in out Unit);
procedure Configure (U : in out Unit; Pin : ESP32S3.GPIO.Pin_Id;
                     Both_Edges : Boolean := False);
function  Count  (U : Unit) return Integer;   --  signed, wraps at +/- 32768
procedure Clear  (U : Unit);
procedure Pause  (U : Unit);
procedure Resume (U : Unit);
\end{lstlisting}

\begin{lstlisting}
declare
   U : Unit;
begin
   Claim (U, 0);
   Configure (U, Pin => 4);             --  count rising edges on GPIO4
   for I in 1 .. 100 loop               --  the self-test drives the same pad ...
      ESP32S3.GPIO.Set (4);   delay until Clock + Microseconds (20);
      ESP32S3.GPIO.Clear (4); delay until Clock + Microseconds (20);
   end loop;
   pragma Assert (Count (U) = 100);     --  ... and counts the edges back
end;   --  U finalizes: paused and released
\end{lstlisting}

\noindent\textbf{Example:} \texttt{examples/esp32s3\_pcnt\_count}.

\section{SDM --- Sigma-Delta Modulator}\index{SDM driver}
\label{app:int:sdm}

The SDM block (part of the GPIO sigma-delta unit, \texttt{GPIO\_SD}) provides
eight 1-bit density-modulated outputs. Each channel emits a high-frequency pulse
stream whose average density is set by a signed 8-bit value; passed through an
external RC low-pass filter it becomes a cheap analog output for LED dimming, a
bias voltage, or simple audio. The Ada driver in \texttt{ESP32S3.SDM} exposes a
channel as a limited, controlled \texttt{Channel} handle.

\texttt{Configure} routes the output to a pin and starts it at 0\,\% density.
\texttt{Carrier\_Hz} requests the pulse-stream (carrier) frequency: the modulator
runs at the APB clock ($\sim$80\,MHz) divided by an integer 1..256, so the
achieved rate is the nearest \texttt{APB/N}, roughly 312\,500\,Hz to 80\,MHz. A
higher carrier is easier to smooth; the exact value rarely matters and is rounded
to the nearest available divider. \texttt{Set\_Density} is a single register
write.

\begin{lstlisting}
type Channel_Index   is range 0 .. 7;      --  the eight SDM channels
subtype Density_Percent is Float range 0.0 .. 100.0;
type Channel is limited private;

procedure Claim     (C : in out Channel; Index : Channel_Index);
function  Is_Valid  (C : Channel) return Boolean;
procedure Release   (C : in out Channel);
procedure Configure (C : in out Channel; Pin : ESP32S3.GPIO.Pin_Id;
                     Carrier_Hz : Positive := 1_000_000);
procedure Set_Density (C : Channel; Percent : Density_Percent);
\end{lstlisting}

\begin{lstlisting}
declare
   Ch        : Channel;
   Densities : constant array (1 .. 3) of Density_Percent := (25.0, 50.0, 75.0);
begin
   Claim (Ch, 0);
   --  A lower carrier oversamples cleanly when the pad is GPIO-sampled back:
   Configure (Ch, Pin => 4, Carrier_Hz => 400_000);
   for I in Densities'Range loop
      Set_Density (Ch, Densities (I));
      delay until Clock + Milliseconds (5);
      --  averaged pad reading should track the set density within a few percent
   end loop;
end;   --  Ch finalizes: output low, channel released
\end{lstlisting}

\noindent\textbf{Example:} \texttt{examples/esp32s3\_sdm\_output}.

\section{RMT --- Remote Control Transceiver}\index{RMT driver}
\label{app:int:rmt}

RMT is the ``infinitely-flexible pulse generator'': it transmits and receives
sequences of \{level, duration\} pulses, making it the workhorse for IR remotes,
WS2812 (``NeoPixel'') LEDs, 1-Wire, and arbitrary bit-banged timing. The S3 has
eight channels --- 0..3 transmit, 4..7 receive --- each backed by a 48-symbol RAM
block. One \texttt{RMT\_Symbol} packs two \{level, duration\} pulses into 32 bits
(durations are 15-bit channel ticks, where the tick length is
$1/\texttt{Resolution\_Hz}$).

The Ada driver in \texttt{ESP32S3.RMT} uses \emph{distinct} limited, controlled
handle types --- \texttt{TX\_Channel} and \texttt{RX\_Channel} --- so transmit and
receive channels can't be confused. \texttt{Transmit} blocks until the burst
finishes; bursts larger than the channel's RAM are streamed by re-filling the
symbol RAM in halves as it drains, so \texttt{Symbols} may be any length. The
\texttt{Blocks} parameter (1..4) borrows the RAM of higher-numbered TX channels
to raise the one-shot ceiling. On the receive side, \texttt{Start} arms the
channel and \texttt{Receive} blocks until the line stays idle for
\texttt{Idle\_Ticks}. Because \texttt{Transmit} busy-polls its re-fill, keep
higher-priority interrupts short.

\begin{lstlisting}
type TX_Index    is range 0 .. 3;   type RX_Index is range 0 .. 3;
type Tick_Count  is range 0 .. 32_767;     --  duration in channel ticks (15-bit)

type RMT_Symbol is record
   Level0 : Boolean := False;   Duration0 : Tick_Count := 0;
   Level1 : Boolean := False;   Duration1 : Tick_Count := 0;
end record;
for RMT_Symbol'Size use 32;
type Symbol_Array is array (Natural range <>) of RMT_Symbol;

type TX_Channel is limited private;   type RX_Channel is limited private;

--  Transmit:
procedure Claim     (C : in out TX_Channel; Index : TX_Index);
procedure Configure (C : in out TX_Channel; Resolution_Hz : Positive;
                     Pin : ESP32S3.GPIO.Pin_Id; Blocks : Positive := 1);
procedure Transmit  (C : TX_Channel; Symbols : Symbol_Array);
--  Receive:
procedure Claim     (C : in out RX_Channel; Index : RX_Index);
procedure Configure (C : in out RX_Channel; Resolution_Hz : Positive;
                     Pin : ESP32S3.GPIO.Pin_Id; Idle_Ticks : Tick_Count := 2_000);
procedure Start     (C : RX_Channel);
procedure Receive   (C : RX_Channel; Into : out Symbol_Array; Count : out Natural);
\end{lstlisting}

\begin{lstlisting}
declare
   Tx   : TX_Channel;   Rx : RX_Channel;
   Sent : constant Symbol_Array :=
     ((Level0 => True, Duration0 =>  50, Level1 => False, Duration1 =>  60),
      (Level0 => True, Duration0 =>  80, Level1 => False, Duration1 =>  90),
      (Level0 => True, Duration0 => 120, Level1 => False, Duration1 => 130));
   Got   : Symbol_Array (0 .. 15);
   Count : Natural;
begin
   Claim (Tx, 0);   Claim (Rx, 0);
   Configure (Tx, Resolution_Hz => 1_000_000, Pin => 4);            --  1 us / tick
   Configure (Rx, Resolution_Hz => 1_000_000, Pin => 4, Idle_Ticks => 1_000);
   Start    (Rx);              --  arm the receiver first
   Transmit (Tx, Sent);        --  drive the burst onto the pad
   Receive  (Rx, Got, Count);  --  block until idle, read it back (single-pad loopback)
end;   --  Tx, Rx finalize: released
\end{lstlisting}

\noindent\textbf{Example:} \texttt{examples/esp32s3\_rmt\_loopback}.

\section{GDMA --- General DMA}\index{GDMA driver}
\label{app:int:gdma}

The ESP32-S3 has a single AHB GDMA block with five channel pairs (\texttt{0..4}).
Each pair carries an independent transmit (OUT) and receive (IN) path, and either
can be wired to any peripheral through the GDMA crossbar --- a channel is a
runtime-assigned resource, not fixed per peripheral. The Ada driver
(\texttt{ESP32S3.GDMA}) models that directly: \texttt{Claim} hands out a free
\texttt{Channel} bound to a \texttt{Peripheral} (\texttt{Mem2Mem}, \texttt{SPI2},
\texttt{I2S0}, \texttt{LCD\_CAM}, \texttt{ADC\_DAC}, \ldots); peripheral drivers
such as I2S and LCD claim a channel and drive their own descriptors over it.

A \texttt{Channel} is \texttt{limited private} --- non-copyable, so two tasks
cannot alias one channel --- and \emph{controlled}: claiming goes through a
protected allocator, and the handle auto-releases on scope exit (including during
exception unwinding) so a channel cannot leak or be reused through a stale copy.
Because finalization is involved, this driver targets the embedded/full profile
(excluded from the light-tasking build). Transfers are bounded by
\texttt{Max\_Transfer = 4095} (the 12-bit buffer-size field); the implemented
memory-to-memory \texttt{Copy} loops one channel's OUT path into its own IN path,
and source, destination and descriptors must all live in internal SRAM. Beyond
\texttt{Copy}, the peripheral-bound \texttt{Start}/\texttt{Start\_Loop}/%
\texttt{Done}/\texttt{Wait} API arms a single buffer (or a self-looping gapless
one) in a \texttt{Direction} and lets the peripheral pace the move.

\begin{lstlisting}
type Channel_Id is mod 5;
type Peripheral is
  (Mem2Mem, SPI2, SPI3, UHCI0, I2S0, I2S1, LCD_CAM, AES, SHA, ADC_DAC, RMT);
type Channel is limited private;            --  non-copyable, auto-released
type Direction is (Mem_To_Periph, Periph_To_Mem);
Max_Transfer : constant := 4095;

procedure Claim   (C : in out Channel; Peri : Peripheral);
function  Is_Valid (C : Channel) return Boolean;
procedure Release (C : in out Channel);
procedure Copy    (C : Channel; Dst, Src : System.Address; Length : Natural);
procedure Start   (C : Channel; Dir : Direction;
                   Buffer : System.Address; Length : Natural);
procedure Start_Loop (C : Channel; Buffer : System.Address; Length : Natural);
function  Done    (C : Channel; Dir : Direction) return Boolean;
procedure Wait    (C : Channel; Dir : Direction);
\end{lstlisting}

\begin{lstlisting}
declare
   C  : Channel;                            --  limited: cannot be aliased
   Ok : Boolean := False;
begin
   Claim (C, Mem2Mem);
   if Is_Valid (C) then
      Copy (C, Dst'Address, Src'Address, Buffer'Length);
      Ok := (for all I in Buffer'Range => Dst (I) = Src (I));
   end if;
   Copy_Result (Ok);
end;                                        --  C finalizes -> channel released
\end{lstlisting}

\cnote{In ESP-IDF a channel is a \texttt{gdma\_channel\_handle\_t} you
\texttt{gdma\_new\_channel} / \texttt{gdma\_del\_channel} by hand, with descriptor
lifetimes left to you. Here the limited-controlled \texttt{Channel} \emph{is} the
allocation: leaving scope frees it, and the type system forbids the aliasing that
leaks one.}

\noindent\textbf{Example:} \texttt{examples/esp32s3\_gdma\_copy}.

\section{I2S --- Digital Audio Serial Bus}\index{I2S driver}
\label{app:int:i2s}

The S3 has two I2S controllers (\texttt{I2S0}, \texttt{I2S1}), each moving PCM
samples only through the GDMA crossbar --- the block has no CPU FIFO, so data
flows exclusively via DMA. The Ada driver (\texttt{ESP32S3.I2S}) brings a port up
as a stereo master that generates BCLK and WS and shifts a sample buffer out (TX)
and/or captures one in (RX), in either \texttt{Standard} (verbatim I2S/TDM) or
\texttt{PDM} mode. In PDM mode the on-chip sigma-delta converters sit between the
buffer and the wire: TX turns each PCM sample into a 1-bit pulse-density stream,
RX decimates a 1-bit PDM input back to PCM, while \texttt{Write}/\texttt{Read}/%
\texttt{Transfer} stay unchanged (the converters high-pass filter, so a constant
level does not survive a PDM round trip).

Mirroring \texttt{ESP32S3.SPI}, the raw register access lives in the private child
\texttt{ESP32S3.I2S.Engine} (which claims the GDMA channel); the application sees
only the parent. \texttt{Setup} configures a port once at startup; thereafter
\texttt{Acquire} hands out a \texttt{Session} --- limited, non-copyable, and
controlled --- that owns the port exclusively and auto-releases on scope exit, so a
fault between \texttt{Acquire} and \texttt{Release} cannot leak the lock.
Transfers (\texttt{Write}, \texttt{Read}, \texttt{Transfer}, plus the gapless
\texttt{Start\_Continuous}/\texttt{Capture}/\texttt{Stop} streaming path) run
outside the lock and raise \texttt{Not\_Owned} if the session does not hold a
port, or \texttt{Not\_Initialized} if a port was never \texttt{Setup}. The driver
needs a tasking runtime (Jorvik or richer) and, via finalization, targets the
embedded/full profile.

\begin{lstlisting}
type I2S_Port    is (I2S0, I2S1);
type Sample_Bits is (Bits_8, Bits_16, Bits_24, Bits_32);
type I2S_Mode    is (Standard, PDM);
type Session     is limited private;        --  non-copyable, auto-released
Not_Initialized, Not_Owned : exception;

procedure Setup (Port : I2S_Port; Sample_Rate : Positive := 16_000;
                 Bits : Sample_Bits := Bits_16; Mode : I2S_Mode := Standard);
procedure Enable_Loopback (Port : I2S_Port; Pad : ESP32S3.GPIO.Pin_Id);
procedure Configure_Pins  (Port : I2S_Port;
                           Bclk, Ws, Dout, Din, Mclk : Optional_Pin := No_Pin);
procedure Acquire  (S : in out Session; Port : I2S_Port);
procedure Write    (S : Session; Tx : System.Address; Length : Natural);
procedure Read     (S : Session; Rx : System.Address; Length : Natural);
procedure Transfer (S : Session; Tx, Rx : System.Address; Length : Natural);
procedure Start_Continuous (S : Session; Tx : System.Address; Length : Natural);
procedure Capture  (S : Session; Rx : System.Address; Length : Natural);
procedure Stop     (S : Session);
procedure Release  (S : in out Session);
\end{lstlisting}

\begin{lstlisting}
Setup (I2S0, Sample_Rate => 16_000, Bits => Bits_16);
Enable_Loopback (I2S0, Pad => Data_Pin);    --  out->in on one pad, no wiring
declare
   S  : Session;                            --  limited: cannot be copied/shared
   Ok : Boolean := False;
begin
   Acquire (S, I2S0);                       --  suspends until the port is free
   Transfer (S, Tx'Address, Rx'Address, Samples'Length * 2);
   Ok := (for all I in Samples'Range => Rx (I) = Tx (I));
end;                                        --  S finalizes -> port released
\end{lstlisting}

\noindent\textbf{Example:} \texttt{examples/esp32s3\_i2s\_loopback} (Standard
full-duplex) and \texttt{examples/esp32s3\_i2s\_pdm} (PDM microphone capture).

\section{LCD --- i8080 Parallel Display}\index{LCD driver}
\label{app:int:lcd}

The LCD half of the S3's LCD\_CAM controller drives an 8-bit Intel-8080
(``i80''/MCU) parallel display: a byte buffer is streamed out the data bus, one
byte per pixel clock (PCLK), over the GDMA crossbar. The Ada driver
(\texttt{ESP32S3.LCD}) covers exactly this 8-bit i80 transmit path (the
camera-receive half and the RGB/16-bit modes are not modelled). \texttt{Setup}
brings the controller up at (about) \texttt{Pclk\_Hz} --- quantised as
\texttt{20\,MHz / round(20\,MHz / Pclk\_Hz)} --- and claims its GDMA channel.

The single controller is guarded by a protected object; \texttt{Acquire} hands out
a limited, controlled \texttt{Session} that owns it exclusively and releases on
scope exit. Notably, both transfers \emph{and} post-\texttt{Setup} reconfiguration
go through the held session: \texttt{Configure\_Pins} (route the eight data lines
\texttt{Data\_Pins} and PCLK) and \texttt{Enable\_Clock\_Out} (free-run the clock
with no data) require ownership, so changing a setting can never race another
task. \texttt{Transmit} streams \texttt{1..4095} bytes from internal SRAM and
reports completion in \texttt{Ok}. Every operation raises \texttt{Not\_Owned}
without the lock, or \texttt{Not\_Initialized} if \texttt{Setup} did not succeed
(including a failed channel claim). All register access lives in the private
\texttt{ESP32S3.LCD.Engine} child; via finalization the driver targets the
embedded/full profile.

\begin{lstlisting}
No_Pin : constant ESP32S3.GPIO.Pad_Number := ESP32S3.GPIO.No_Pin;
type Data_Pins is array (0 .. 7) of ESP32S3.GPIO.Optional_Pin;  -- D0..D7
type Session   is limited private;          --  non-copyable, auto-released
Not_Initialized, Not_Owned : exception;

procedure Setup (Pclk_Hz : Positive := 1_000_000);
procedure Acquire (S : in out Session);
procedure Configure_Pins   (S : Session; Data : Data_Pins;
                            Pclk : ESP32S3.GPIO.Optional_Pin);
procedure Enable_Clock_Out (S : Session; Pclk_Pad : ESP32S3.GPIO.Pin_Id);
procedure Transmit (S : Session; Tx : System.Address; Length : Natural;
                    Ok : out Boolean);
procedure Release  (S : in out Session);
\end{lstlisting}

\begin{lstlisting}
Setup (Pclk_Hz => 200 * 1000);              --  200 kHz pixel clock
declare
   S  : Session;
   Ok : Boolean;
begin
   Acquire (S);
   Configure_Pins (S, D_Pins, Pclk => Pclk_Pin);  --  route pads on held ctrl
   Transmit (S, Buf'Address, Buffer'Length, Ok);  --  blocks until done
   Release (S);
   Put_Line (if Ok then "PASS" else "FAIL");
end;
\end{lstlisting}

\noindent\textbf{Example:} \texttt{examples/esp32s3\_lcd\_i8080}.

\section{ADC --- SAR Analog-to-Digital Converter}\index{ADC driver}
\label{app:int:adc}

The S3 has two SAR ADC units, each with up to ten 12-bit channels on fixed GPIOs:
ADC1 channel \texttt{n} maps to \texttt{GPIO(n+1)} (ch0\,=\,GPIO1\ldots ch9\,=\,GPIO10),
ADC2 channel \texttt{n} to \texttt{GPIO(n+11)}. The Ada driver
(\texttt{ESP32S3.ADC}) does software-triggered single (one-shot) conversions
through the RTC controller; each returns a raw \texttt{0..4095} code, with
per-channel \texttt{Attenuation} (\texttt{Db\_0} $\approx$ 1.1\,V full-scale
\ldots \texttt{Db\_12} $\approx$ 3.3\,V) setting the input range.

A unit is a shared resource handed out as a \texttt{Reader} --- limited
(non-copyable) and controlled (released automatically on scope exit) --- so
claiming a unit cannot leak. After \texttt{Claim}, check \texttt{Is\_Valid};
\texttt{Read} samples a channel (blocking, a few microseconds), returning
\texttt{0} on an invalid handle. \texttt{Channel\_Pin} reports the GPIO a channel
is wired to (routing/documentation), and the diagnostics \texttt{Cal\_Code}
(self-calibrated initial code) and \texttt{Last\_Done} (the SAR DONE flag) help
confirm a conversion actually ran. Via finalization the driver targets the
embedded/full profile.

\begin{lstlisting}
type ADC_Unit      is (ADC1, ADC2);
type Channel_Index is range 0 .. 9;
type Attenuation   is (Db_0, Db_2_5, Db_6, Db_12);
subtype Raw_Value  is Natural range 0 .. 4095;        --  12-bit result
type Reader        is limited private;                --  auto-released

procedure Claim   (R : in out Reader; Unit : ADC_Unit);
function  Is_Valid (R : Reader) return Boolean;
procedure Release (R : in out Reader);
function  Read    (R : Reader; Ch : Channel_Index;
                   Atten : Attenuation := Db_12) return Raw_Value;
function  Channel_Pin (Unit : ADC_Unit; Ch : Channel_Index)
                       return ESP32S3.GPIO.Pin_Id;
function  Cal_Code  (Unit : ADC_Unit) return Natural;
function  Last_Done return Boolean;
\end{lstlisting}

\begin{lstlisting}
declare
   R : Reader;
   High, Low : Natural := 0;
begin
   Claim (R, ADC1);                         --  ADC1 ch0 -> GPIO1
   ESP32S3.GPIO.Configure (Pin, ESP32S3.GPIO.Output);
   ESP32S3.GPIO.Set (Pin);                  --  drive the sensed pad high
   delay until Clock + Milliseconds (2);
   High := Read (R, Ch);                    --  expect near full scale
   ESP32S3.GPIO.Clear (Pin);
   delay until Clock + Milliseconds (2);
   Low := Read (R, Ch);                     --  expect near zero
end;                                        --  R finalizes -> unit released
\end{lstlisting}

\noindent\textbf{Example:} \texttt{examples/esp32s3\_adc\_read}.

\section{TWAI --- CAN Controller}\index{TWAI driver}
\label{app:int:twai}

The S3's TWAI block (Two-Wire Automotive Interface) is a single
SJA1000-compatible CAN 2.0 controller. The Ada driver (\texttt{ESP32S3.TWAI})
sends and receives both standard (11-bit \texttt{Standard\_Id}) and extended
(29-bit \texttt{Extended\_Id}, CAN 2.0B) frames, data or remote-request (RTR).
Each frame type carries its own identifier subtype, so the id is range-checked to
its standard and \texttt{Send} is overloaded on the frame type --- a 29-bit id
cannot end up in a standard frame. A real bus needs an external transceiver; for a
wiring-free self-test, \texttt{Self\_Test} mode transmits and self-receives a
frame (no second node to acknowledge), with TX looped back to RX through one pad
via \texttt{Enable\_Loopback}.

The single controller is guarded by a protected object; \texttt{Acquire} hands out
a limited, controlled \texttt{Session} owning it exclusively, auto-released on
scope exit. \texttt{Setup} (mode + bit rate, accepting all ids) runs once;
transfers and post-\texttt{Setup} reconfiguration (\texttt{Configure\_Pins},
\texttt{Enable\_Loopback}) all go through the held session and raise
\texttt{Not\_Owned} otherwise (or \texttt{Not\_Initialized} before \texttt{Setup}).
Because a received frame may be either width, receiving is a two-step idiom:
\texttt{Available} reports a waiting frame, \texttt{Is\_Extended} its width, then
the matching \texttt{Receive} overload reads it (returning \texttt{Got => False}
if the waiting frame is the other width, leaving it buffered). All register access
lives in the private \texttt{ESP32S3.TWAI.Engine} child; via finalization the
driver targets the embedded/full profile.

\begin{lstlisting}
type Bus_Mode    is (Normal, Listen_Only, Self_Test);
subtype Data_Length is Natural range 0 .. 8;
type Data_Bytes  is array (0 .. 7) of Interfaces.Unsigned_8;
subtype Standard_Id is Interfaces.Unsigned_32 range 0 .. 16#7FF#;        -- 11-bit
subtype Extended_Id is Interfaces.Unsigned_32 range 0 .. 16#1FFF_FFFF#;  -- 29-bit
type Standard_Frame is record
   Id : Standard_Id := 0; Remote : Boolean := False;
   Length : Data_Length := 0; Data : Data_Bytes := (others => 0);
end record;
type Extended_Frame is record ... end record;   --  same shape, Extended_Id
type Session is limited private;

procedure Setup (Mode : Bus_Mode := Normal; Bit_Rate : Positive := 125_000);
procedure Acquire (S : in out Session);
procedure Enable_Loopback (S : Session; Pad : ESP32S3.GPIO.Pin_Id);
procedure Send (S : Session; F : Standard_Frame);   --  overloaded on frame type
procedure Send (S : Session; F : Extended_Frame);
function  Available   (S : Session) return Boolean;
function  Is_Extended (S : Session) return Boolean;
procedure Receive (S : Session; F : out Standard_Frame; Got : out Boolean);
procedure Receive (S : Session; F : out Extended_Frame; Got : out Boolean);
\end{lstlisting}

\begin{lstlisting}
Setup (Mode => Self_Test, Bit_Rate => 125_000);
declare
   S   : Session;
   RS  : Standard_Frame;
   Got : Boolean;
begin
   Acquire (S);
   Enable_Loopback (S, Pad);                --  out->in on one pad, no wiring
   Send (S, Std);                           --  picks the Standard_Frame overload
   Got := Available (S) and then not Is_Extended (S);
   if Got then
      Receive (S, RS, Got);                 --  matching width -> Got => True
   end if;
end;                                        --  S finalizes -> controller released
\end{lstlisting}

\noindent\textbf{Example:} \texttt{examples/esp32s3\_twai\_loopback}.

\section{RTC --- Real-Time Clock and Deep Sleep}\index{RTC driver}
\label{app:int:rtc}

The ESP32-S3's RTC (low-power) domain stays alive while the digital core is
powered down. In deep sleep, CPU and main RAM lose power; only the RTC domain
survives, so on wake the chip \emph{resets} and re-runs from the start --- but
anything kept in the 8\,KB of RTC slow memory (at \texttt{0x5000\_0000}) persists
across the reset. \texttt{ESP32S3.RTC} exposes that retained region, the cause of
the current boot (power-on versus a deep-sleep wake), and deep-sleep entry with a
timer or an RTC-GPIO (EXT1) wake source.

The deep-sleep timer runs on the uncalibrated RTC slow clock
(\textasciitilde136\,kHz), so \texttt{Deep\_Sleep\_For} is approximate. The entry
procedures do not return on success: the chip powers its digital core down and
resets on wake. Retained memory can be reached word-by-word with
\texttt{Read}/\texttt{Write}, by overlaying a variable on the \texttt{Slow\_Memory}
address, or through the generic \texttt{Retained} package. The operations are plain
register pokes with no tasking or finalization, so the driver works under every
runtime profile.

\begin{lstlisting}
Slow_Memory      : constant System.Address := ...;   --  0x5000_0000
Slow_Memory_Size : constant := 8 * 1024;
subtype Word_Index is Natural range 0 .. Slow_Memory_Size / 4 - 1;

function  Read  (Index : Word_Index) return Interfaces.Unsigned_32;
procedure Write (Index : Word_Index; Value : Interfaces.Unsigned_32);

generic                                  --  typed retained object
   type Item is private;
   Offset : Natural := 0;
package Retained is
   Object : Item with Import, Volatile, Address => ...;
end Retained;

type Wake_Cause is (Power_On, Deep_Sleep_Timer, Deep_Sleep_GPIO, Other_Reset);
function Last_Wake return Wake_Cause;

procedure Deep_Sleep_For   (Wake_After : Duration);            --  timer wake
procedure Deep_Sleep_Until (Pin : ESP32S3.GPIO.Pin_Id;        --  EXT1 wake
                            High : Boolean := True);
procedure Disable_Super_Watchdog;
\end{lstlisting}

A boot counter kept in retained word 0 survives the deep-sleep reset; the wake
cause turns into \texttt{Deep\_Sleep\_Timer} on each timer wake.

\begin{lstlisting}
use ESP32S3.RTC;
WC         : constant Wake_Cause := Last_Wake;
Boot_Count : Unsigned_32 := Read (0);          --  retained across deep sleep
begin
   Disable_Super_Watchdog;                     --  a wake can leave it armed
   if WC = Deep_Sleep_Timer or else WC = Deep_Sleep_GPIO then
      Boot_Count := Boot_Count + 1;
   else
      Boot_Count := 1;                          --  power-on / flash reset
   end if;
   Write (0, Boot_Count);                       --  persist it

   if Boot_Count < 4 then
      Deep_Sleep_For (2.0);                      --  ~2 s; DOES NOT RETURN
   end if;                                       --  chip resets on wake
\end{lstlisting}

\noindent\textbf{Example:} \texttt{examples/esp32s3\_rtc\_sleep}.

\section{RTC-IO --- Low-Power GPIO and Pad Hold}\index{RTC\_IO driver}
\label{app:int:rtcio}

GPIO0 through GPIO21 are RTC-capable: they belong to the low-power IO domain. The
headline feature is pad \emph{hold} --- latch a pad at its current output level so
it keeps driving while the rest of the chip changes, and, crucially, while the
chip is in deep sleep (the digital core powers down, but a held RTC pad stays put)
and across the reset a deep-sleep wake causes. Use it to keep a load enabled or a
reset line asserted while you sleep. A held pad ignores ordinary GPIO writes until
you \texttt{Release} it.

\texttt{ESP32S3.RTC\_IO} also exposes the RTC-domain pulls, which are distinct
from the digital IO\_MUX pulls configured through \texttt{ESP32S3.GPIO}: route a
pad into the RTC domain with \texttt{Enable\_RTC\_Input}, then \texttt{Set\_Pull}
gives it a defined level that holds (including in deep sleep). Register pokes only,
so RTC-IO works under every runtime profile.

\begin{lstlisting}
subtype RTC_Pin is ESP32S3.GPIO.Pin_Id range 0 .. 21;

procedure Hold    (Pin : RTC_Pin);    --  latch pad at its current level
procedure Release (Pin : RTC_Pin);    --  follow the GPIO register again
function  Is_Held (Pin : RTC_Pin) return Boolean;

procedure Enable_RTC_Input (Pin : RTC_Pin);   --  route into RTC domain
type Pull_Mode is (No_Pull, Up, Down);
procedure Set_Pull (Pin : RTC_Pin; Mode : Pull_Mode);
\end{lstlisting}

Driving a pad high and holding it makes a subsequent GPIO clear a no-op until
release:

\begin{lstlisting}
Pin : constant ESP32S3.GPIO.Pin_Id := 5;     --  RTC-capable pad
begin
   ESP32S3.GPIO.Configure (Pin, ESP32S3.GPIO.Output);
   ESP32S3.GPIO.Set (Pin);                    --  drive high

   Hold (Pin);                                --  latch it high
   ESP32S3.GPIO.Clear (Pin);                  --  IGNORED while held
   --  Read (Pin) still reports high here

   Release (Pin);                             --  unlatch
   ESP32S3.GPIO.Clear (Pin);                  --  now takes effect -> low
\end{lstlisting}

\noindent\textbf{Example:} \texttt{examples/esp32s3\_rtcio\_hold}.

\section{Touch --- Capacitive Touch Sensor}\index{Touch driver}
\label{app:int:touch}

The S3's capacitive touch sensor (v2) provides 14 channels, channel $n$ wired to
GPIO$n$ ($n = 1\ldots14$). Each channel measures the self-capacitance of its pad
by counting charge/discharge cycles in a fixed window; a finger near the pad
raises the capacitance and changes the count. A hardware FSM scans the enabled
channels continuously on the RTC timer, and \texttt{Read} returns the latest raw
count.

\texttt{ESP32S3.Touch} keeps the model deliberately simple: bring the controller
up once with \texttt{Setup}, \texttt{Enable} the channels you want, then
\texttt{Read} their raw counts. \texttt{Touched} is a software comparison against a
caller-captured \texttt{Reference} with a \texttt{Margin} --- deterministic, with
no on-chip threshold registers in play. It lives in the RTC/SENS domain and is
register pokes only, so it works under every runtime profile.

\begin{lstlisting}
type Channel is range 1 .. 14;               --  channel n -> GPIO n
function Pad (Ch : Channel) return ESP32S3.GPIO.Pin_Id;

procedure Setup;                             --  start the scanning FSM (once)
procedure Enable (Ch : Channel);            --  add a channel to the scan
function  Read   (Ch : Channel) return Natural;   --  latest raw count

function Touched (Ch        : Channel;
                  Reference : Natural;
                  Margin    : Natural := 20_000) return Boolean;
\end{lstlisting}

A bare pad reads a stable non-zero baseline; two different pads read different
values, proving the FSM runs on silicon:

\begin{lstlisting}
A : constant Channel := 1;     --  GPIO1
B : constant Channel := 3;     --  GPIO3
begin
   Setup;
   Enable (A);
   Enable (B);
   delay until Clock + Milliseconds (50);     --  let the FSM scan

   declare
      Ra : constant Natural := Read (A);
      Rb : constant Natural := Read (B);
   begin
      --  Ra, Rb both non-zero and distinct => capacitance FSM is live
      null;
   end;

   --  Touch detection against a captured baseline:
   Baseline : constant Natural := Read (A);
   Hit      : constant Boolean := Touched (A, Baseline, Margin => 50_000);
\end{lstlisting}

\noindent\textbf{Example:} \texttt{examples/esp32s3\_touch\_read}.

\section{SHA --- Hardware Hash Accelerator}\index{SHA driver}
\label{app:int:sha}

The SHA accelerator computes SHA-1, SHA-224 and SHA-256 of a byte message in
hardware. All three variants share the same 512-bit block and padding, differing
only in the hardware mode and digest length. \texttt{ESP32S3.SHA} exposes them as
three one-shot functions that take a \texttt{Byte\_Array} of any length (padded
internally) and return the fixed-size digest.

The accelerator is a single shared resource, so a protected object serializes the
message-load / start / read handshake, making concurrent \texttt{Hash} calls from
different tasks safe. There is no finalization --- just register pokes --- so the
driver works under every runtime profile. (The block hardware also does
SHA-384/512 on a 1024-bit block; those are not exposed here.)

\begin{lstlisting}
type Byte_Array is array (Natural range <>) of Interfaces.Unsigned_8;

subtype SHA1_Digest   is Byte_Array (0 .. 19);   --  160-bit
subtype SHA224_Digest is Byte_Array (0 .. 27);   --  224-bit
subtype SHA256_Digest is Byte_Array (0 .. 31);   --  256-bit

function Hash_1   (Data : Byte_Array) return SHA1_Digest;
function Hash_224 (Data : Byte_Array) return SHA224_Digest;
function Hash_256 (Data : Byte_Array) return SHA256_Digest;
\end{lstlisting}

Hashing \texttt{"abc"} and comparing against the FIPS-180 vector:

\begin{lstlisting}
use type ESP32S3.SHA.Byte_Array;

Abc : constant ESP32S3.SHA.Byte_Array := (16#61#, 16#62#, 16#63#);  --  "abc"

Sha_Expect : constant ESP32S3.SHA.SHA256_Digest :=
  (16#BA#, 16#78#, 16#16#, 16#BF#, 16#8F#, 16#01#, 16#CF#, 16#EA#,
   16#41#, 16#41#, 16#40#, 16#DE#, 16#5D#, 16#AE#, 16#22#, 16#23#,
   16#B0#, 16#03#, 16#61#, 16#A3#, 16#96#, 16#17#, 16#7A#, 16#9C#,
   16#B4#, 16#10#, 16#FF#, 16#61#, 16#F2#, 16#00#, 16#15#, 16#AD#);
begin
   --  PASS when the hardware digest matches the published vector:
   if ESP32S3.SHA.Hash_256 (Abc) = Sha_Expect then ... end if;
\end{lstlisting}

\noindent\textbf{Example:} \texttt{examples/esp32s3\_crypto}.

\section{AES --- Hardware AES Accelerator}\index{AES driver}
\label{app:int:aes}

The AES accelerator performs single-block (128-bit) AES in hardware.
\texttt{ESP32S3.AES} exposes the ECB building block --- one-block encrypt and
decrypt --- for 128- and 256-bit keys. As with SHA, the engine is a single shared
resource guarded by a protected object, so calls from different tasks are safe,
and the register-poke implementation needs no finalization and works under every
runtime profile.

\cnote{The ESP32-S3 AES hardware supports only 128- and 256-bit keys; there is no
192-bit mode on this silicon (selecting ``AES-192'' silently falls back to AES-128
on the first 16 key bytes). The \texttt{Supported\_Key} precondition turns an
unsupported key length into a contract violation instead of a silent fallback ---
callers that pass the \texttt{Key\_128} / \texttt{Key\_256} subtypes satisfy it
statically, at no run-time cost.}

\begin{lstlisting}
type Block is array (0 .. 15) of Interfaces.Unsigned_8;   --  128-bit block

type Key_Bytes is array (Natural range <>) of Interfaces.Unsigned_8;
subtype Key_128 is Key_Bytes (0 .. 15);
subtype Key_256 is Key_Bytes (0 .. 31);

function Supported_Key (Key : Key_Bytes) return Boolean is
  (Key'Length = 16 or else Key'Length = 32);

function Encrypt_ECB (Key : Key_Bytes; Plain  : Block) return Block
  with Pre => Supported_Key (Key);
function Decrypt_ECB (Key : Key_Bytes; Cipher : Block) return Block
  with Pre => Supported_Key (Key);
\end{lstlisting}

Encrypting the FIPS-197 example block and round-tripping it back:

\begin{lstlisting}
use type ESP32S3.AES.Block;

Aes_Key : constant ESP32S3.AES.Key_128 :=
  (16#00#, 16#01#, 16#02#, 16#03#, 16#04#, 16#05#, 16#06#, 16#07#,
   16#08#, 16#09#, 16#0A#, 16#0B#, 16#0C#, 16#0D#, 16#0E#, 16#0F#);
Aes_Plain : constant ESP32S3.AES.Block :=
  (16#00#, 16#11#, 16#22#, 16#33#, 16#44#, 16#55#, 16#66#, 16#77#,
   16#88#, 16#99#, 16#AA#, 16#BB#, 16#CC#, 16#DD#, 16#EE#, 16#FF#);
Aes_Cipher : constant ESP32S3.AES.Block :=
  (16#69#, 16#C4#, 16#E0#, 16#D8#, 16#6A#, 16#7B#, 16#04#, 16#30#,
   16#D8#, 16#CD#, 16#B7#, 16#80#, 16#70#, 16#B4#, 16#C5#, 16#5A#);
begin
   --  encrypt matches the vector, and decrypt recovers the plaintext:
   if ESP32S3.AES.Encrypt_ECB (Aes_Key, Aes_Plain)  = Aes_Cipher and then
      ESP32S3.AES.Decrypt_ECB (Aes_Key, Aes_Cipher) = Aes_Plain
   then ... end if;
\end{lstlisting}

\noindent\textbf{Example:} \texttt{examples/esp32s3\_crypto}.

\section{SDMMC --- Native SD/MMC Host Controller}\index{SDMMC driver}
\label{app:int:sdmmc}

The dedicated SDHOST controller drives an SD card on the real SD bus --- a clock,
a bidirectional command line, and 1 or 4 data lines --- with its own
command/response state machine. Unlike the SPI-mode SD path, this is how you reach
an SDHC/SDXC card at speed (up to 50\,MHz in High-Speed mode).
\texttt{ESP32S3.SDMMC} moves data in PIO/FIFO mode (the CPU pushes/pops each
512-byte block through the controller FIFO), so there are no GDMA descriptors and
no finalization: a library-level protected object serializes the single shared
controller, and it works under every runtime profile, including light tasking. The
API is always 512-byte logical blocks (LBA); the SDHC/SDXC versus SDSC addressing
difference is resolved internally. (Covered in full in the Mass Storage chapter,
\S\ref{ch:storage}.)

\begin{lstlisting}
type Slot      is (Slot1, Slot2);
type Bus_Width is (Width_1, Width_4);
type Block         is array (0 .. 511) of Interfaces.Unsigned_8;
type Block_Address is new Interfaces.Unsigned_32;
type Card_Kind is (Unknown, SDSC, SDHC);             --  SDHC also covers SDXC
type Status    is (OK, No_Card, Unusable, Init_Timeout, Cmd_Timeout,
                   Cmd_CRC, Read_Error, Write_Error);
type Card is limited private;

procedure Setup (C : out Card; On : Slot;
                 Clk, Cmd, D0 : ESP32S3.GPIO.Pin_Id;
                 D1, D2, D3   : ESP32S3.GPIO.Optional_Pin := ESP32S3.GPIO.No_Pin;
                 Width        : Bus_Width := Width_1;
                 Init_Clock_Hz, Data_Clock_Hz : Positive;
                 High_Speed   : Boolean := False);
procedure Initialize  (C : in out Card; Result : out Status);
function  Kind        (C : Card) return Card_Kind;
procedure Read_Block  (C : in out Card; LBA : Block_Address;
                       Data : out Block; Result : out Status);
procedure Write_Block (C : in out Card; LBA : Block_Address;
                       Data : Block; Result : out Status);
\end{lstlisting}

Bringing up Slot1 in 4-bit mode and doing a non-destructive sector round-trip:

\begin{lstlisting}
use type ESP32S3.SDMMC.Status;
Test_LBA : constant ESP32S3.SDMMC.Block_Address := 16#2000#;   --  sector 8192
C    : ESP32S3.SDMMC.Card;
St   : ESP32S3.SDMMC.Status;
Orig : ESP32S3.SDMMC.Block;
begin
   ESP32S3.SDMMC.Setup
     (C, On => ESP32S3.SDMMC.Slot1,
      Clk => 14, Cmd => 15, D0 => 2, D1 => 4, D2 => 12, D3 => 13,
      Width => ESP32S3.SDMMC.Width_4,
      Init_Clock_Hz => 400_000, Data_Clock_Hz => 20_000_000);

   ESP32S3.SDMMC.Initialize (C, St);
   if St = ESP32S3.SDMMC.OK then
      ESP32S3.SDMMC.Read_Block  (C, Test_LBA, Orig, St);   --  read scratch sector
      ESP32S3.SDMMC.Write_Block (C, Test_LBA, Orig, St);   --  write it back unchanged
   end if;
\end{lstlisting}

\noindent\textbf{Example:} \texttt{examples/esp32s3\_sdmmc}.
